\documentclass[11pt]{article}
\usepackage{amsmath,amssymb,amsthm}
\usepackage{booktabs}
\usepackage{geometry}
\usepackage{algorithm,algpseudocode}
\geometry{margin=1in}

\newcommand{\R}{\mathbb{R}}
\newcommand{\ip}[2]{\langle #1, #2 \rangle}
\newcommand{\norm}[1]{\|#1\|}
\newcommand{\half}{\tfrac{1}{2}}
\DeclareMathOperator{\tr}{tr}
\DeclareMathOperator{\Exp}{Exp}
\DeclareMathOperator{\diag}{diag}

\title{Geodesic Interior-Point Method for Cones\\with Log-Homogeneous Barriers}
\author{}
\date{}

\begin{document}
\maketitle

\section{Problem}

We consider the conic program
\begin{equation}\label{eq:problem}
  \min_{x \in \R^p} \; c^\top x \quad \text{s.t.} \quad Ax + b \in K,
\end{equation}
where $K \subset \R^m$ is a proper cone with a $\nu$-logarithmically homogeneous
self-concordant barrier (LHSCB) $F : \operatorname{int}(K) \to \R$:
\[
  F(t \cdot s) = F(s) - \nu \log t, \quad t > 0.
\]
This implies the scaling identities
\begin{equation}\label{eq:log-homog}
  \nabla F(ts) = \frac{1}{t} \nabla F(s), \qquad
  \nabla^2 F(ts) = \frac{1}{t^2} \nabla^2 F(s).
\end{equation}

\section{Primal-dual variables}

The slack is $s = Ax + b \in \operatorname{int}(K)$.
The dual variable is defined via the barrier gradient:
\[
  \lambda = -\mu \nabla F(s), \qquad \mu = \frac{1}{k^2},
\]
which satisfies $\lambda \in \operatorname{int}(K^*)$ (since $-\nabla F$ maps
$\operatorname{int}(K) \to \operatorname{int}(K^*)$ for any LHSCB).

By Euler's theorem for $\nu$-log-homogeneous functions,
$\ip{s}{\nabla F(s)} = -\nu$, so the complementarity is
\[
  \ip{s}{\lambda} = \mu \nu = \frac{\nu}{k^2}.
\]

\section{KKT system}

At parameter $k$ (equivalently $\mu = 1/k^2$), the primal-dual Newton system is:
\begin{equation}\label{eq:kkt}
  \begin{bmatrix} \mu H(\lambda) & A \\ A^\top & 0 \end{bmatrix}
  \begin{bmatrix} \delta\lambda \\ x \end{bmatrix}
  =
  \begin{bmatrix} s - b \\ c - A^\top \lambda \end{bmatrix},
\end{equation}
where $H(\lambda) = \nabla^2 F^*(\lambda)$ is the dual barrier Hessian.
The variable $x$ is the \emph{full} primal (not an increment), and $\delta\lambda$
is the dual update.  After solving:
\begin{itemize}
  \item $s + \delta s = Ax + b$ \quad (primal feasible by construction),
  \item $A^\top(\lambda + \delta\lambda) = c$ \quad (dual feasible by construction),
  \item $\text{gap} = \ip{s + \delta s}{\lambda + \delta\lambda}$.
\end{itemize}

\subsection{Schur complement}

Using the identity $H(\lambda)^{-1} = \mu^2 \nabla^2 F(s)$ (which holds when
$\lambda = -\mu \nabla F(s)$), the Schur complement gives:
\begin{equation}\label{eq:schur}
  \bigl(A^\top \nabla^2 F(s) \, A\bigr) \, x
  = A^\top\bigl(\nabla^2 F(s)(s - b) - \nabla F(s)\bigr) - k^2 c.
\end{equation}

\subsection{Scaled formulation: $w = k \cdot s$, $z = k \cdot x$}

Substituting $s = w/k$ and using~\eqref{eq:log-homog}:
\[
  \nabla^2 F(s) = k^2 \nabla^2 F(w), \quad \nabla F(s) = k \nabla F(w).
\]
The Schur complement~\eqref{eq:schur} becomes, after dividing by $k$:
\begin{equation}\label{eq:schur-w}
  \bigl(A^\top \nabla^2 F(w) \, A\bigr) \, z
  = A^\top\bigl(\nabla^2 F(w) \, w - \nabla F(w)\bigr)
    - k \bigl(c + A^\top \nabla^2 F(w) \, b\bigr),
\end{equation}
where $z = kx$.  The Gram matrix $A^\top \nabla^2 F(w) A$ is \textbf{independent of $k$}.

\subsection{Two-solve decomposition}

Factor the Gram once, then two back-solves:
\begin{align}
  \bigl(A^\top \nabla^2 F(w) A\bigr) \, z_0 &= A^\top\bigl(\nabla^2 F(w) w - \nabla F(w)\bigr)
    && \text{(centering)}, \label{eq:z0} \\
  \bigl(A^\top \nabla^2 F(w) A\bigr) \, z_1 &= -\bigl(c + A^\top \nabla^2 F(w) b\bigr)
    && \text{(optimality)}. \label{eq:z1}
\end{align}
The combined solution is $z(k) = z_0 + k \cdot z_1$, i.e.,
\[
  x(k) = \frac{z(k)}{k} = \frac{z_0}{k} + z_1.
\]
Both $c$ and $b$ appear in~\eqref{eq:z1}, scaled by $k$ symmetrically.

\section{Line search for $k$}

Define the cone-space directions $d_0 = A z_0$ and $d_1 = A z_1$.  The
displacement in $w$-space is $\delta w(k) = d_0 + k \cdot d_1$.

\paragraph{Dikin bound.}
The largest $k$ with $\norm{\delta w(k)}_{H(w)} \leq 1$ is the root of the
quadratic
\[
  \ip{d_1}{d_1}_H k^2 + 2\ip{d_0}{d_1}_H k + (\ip{d_0}{d_0}_H - 1) = 0,
\]
where $\ip{a}{b}_H = a^\top \nabla^2 F(w) \, b$.

\paragraph{Primal-dual feasibility.}
Beyond the Dikin bound, binary search for the largest $k$ such that:
\begin{enumerate}
  \item $s_{\text{new}} = Ax(k) + b \in \operatorname{int}(K)$ \quad (primal feasibility),
  \item $\lambda_{\text{new}} = \lambda + \delta\lambda \in \operatorname{int}(K^*)$ \quad (dual feasibility).
\end{enumerate}
For low-dimensional cones (e.g., the 3D exponential cone), each feasibility
check is $O(1)$.  The geodesic step preserves feasibility beyond the Dikin
ellipsoid.

\section{Geodesic step}

After selecting $k$, take a geodesic step on $\operatorname{int}(K)$ with
respect to the Riemannian metric $g = \nabla^2 F$:
\[
  s_{\text{new}} = \Exp_s(\alpha \cdot \delta s),
\]
where $\alpha = 1/(1 + \norm{\delta s}_H)$ is the step size for self-concordant
barriers, and $\delta s = s(k) - s$ (the Newton displacement).

The geodesic $\Exp_s$ satisfies the ODE
\[
  \ddot\gamma + \half \nabla^2 F(\gamma)^{-1} T(\dot\gamma, \dot\gamma) = 0,
  \quad T(\dot\gamma, \dot\gamma)_\ell = \dot\gamma^\top \frac{\partial \nabla^2 F}{\partial \gamma_\ell} \dot\gamma.
\]
For low-dimensional cones (e.g., 3D exponential), this is integrated by
St\"ormer--Verlet with $\sim$8 steps.

For symmetric cones, the geodesic reduces to the automorphism update
$W \mapsto W \cdot \exp(\alpha \cdot d)$.

\section{Algorithm}

\begin{algorithm}[H]
\caption{Geodesic LP for cones with log-homogeneous barriers}
\begin{algorithmic}[1]
\State \textbf{Input:} $A, b, c$; barrier $F$ with parameter $\nu$; tolerance $\varepsilon$.
\State Initialize $x = 0$, $k = 1$.
\Repeat
  \State $s \gets Ax + b$, \quad $w \gets k \cdot s$.
  \State Compute $\nabla F(w)$, $\nabla^2 F(w)$.
  \State Factor $G = A^\top \nabla^2 F(w) \, A$.
  \State Solve $G \, z_0 = A^\top(\nabla^2 F(w) \, w - \nabla F(w))$.
    \hfill (centering)
  \State Solve $G \, z_1 = -(c + A^\top \nabla^2 F(w) \, b)$.
    \hfill (optimality)
  \State Line search: find largest $k$ with
    $s(k) \in K$, $\lambda(k) \in K^*$.
  \State $z \gets z_0 + k \cdot z_1$, \quad $x \gets z / k$.
  \State $s_{\text{new}} \gets Ax + b$.
  \State Compute $\lambda_{\text{new}}$ from KKT~\eqref{eq:kkt}, verify
    $A^\top \lambda_{\text{new}} = c$.
  \State $\text{gap} \gets \ip{s_{\text{new}}}{\lambda_{\text{new}}}$.
  \State Geodesic step: $s \gets \Exp_s(\alpha \cdot (s_{\text{new}} - s))$.
  \State Recover $x$ from $s = Ax + b$.
\Until{$\text{gap} < \varepsilon$}
\end{algorithmic}
\end{algorithm}

\section{Dictionary: general cone vs.\ symmetric cone}

The following table maps each object in the general geodesic IPM to its
symmetric cone specialization.

\begin{table}[h]
\centering
\renewcommand{\arraystretch}{1.3}
\begin{tabular}{@{}lll@{}}
\toprule
\textbf{Object} & \textbf{General cone} & \textbf{Symmetric cone} \\
\midrule
Cone $K$ & Any proper cone with LHSCB & $\R^n_+$, $\mathcal{L}^n$ (SOC), $\mathbb{S}^n_+$ (PSD) \\
Barrier $F(s)$ & $\nu$-LHSCB on $K$ & $-\log\det(s)$ \\
Barrier parameter $\nu$ & Depends on $F$ & $n$ (nonneg), $2$ (SOC), $n$ (PSD) \\
Hessian $\nabla^2 F(s)$ & General $m \times m$ SPD matrix & $P(s^{-1})$ (quadratic rep.) \\
Gradient $\nabla F(s)$ & General $m$-vector & $-s^{-1}$ (Jordan inverse) \\
\midrule
Scaling $w$ & $w = k \cdot s$ & $w = k \cdot s$ \\
Gram matrix & $A^\top \nabla^2 F(w) \, A$ & $A^\top P(w^{-1}) \, A$ \\
Centering RHS & $A^\top(\nabla^2 F(w)w - \nabla F(w))$ & $2A^\top w^{-1}$ \\
Optimality RHS & $-(c + A^\top \nabla^2 F(w) b)$ & $-(c + A^\top P(w^{-1}) b)$ \\
\midrule
Dual variable & $\lambda = -\mu \nabla F(s) \in K^*$ & $\lambda = \mu s^{-1} \in K$ \\
Complementarity & $\ip{s}{\lambda} = \nu / k^2$ & $\ip{s}{\lambda} = \nu / k^2$ \\
Dual feasibility & $A^\top \lambda = c$ (from KKT) & $A^\top \lambda = c$ (from KKT) \\
\midrule
Direction $d$ & $d_0 + k \cdot d_1$ in $w$-space & $d_0 + k \cdot d_1$ in $w$-space \\
Line search & Dikin $\|d\|_H \leq 1$ + feasibility & $\|d\|_\infty \leq 1$ (eigenvalue bound) \\
Step size & $\alpha = 1/(1 + \|d\|_H)$ & $\alpha = \min(1, 2/\|d\|_\infty^2)$ \\
\midrule
Geodesic step & $\Exp_s(\alpha \cdot d)$ via ODE & $w \gets w \cdot \exp(-\alpha \cdot d)$ \\
Automorphism & None (general) & $M \in GL(n)$: $s \mapsto MsM^\top$ \\
Update & Recompute $\nabla^2 F$ at new $s$ & $M \gets M \cdot \exp(\alpha d / 2)$ \\
\bottomrule
\end{tabular}
\caption{Dictionary between the general geodesic IPM and the symmetric cone case.
The symmetric cone objects specialize via the Jordan algebra structure:
$P(a)b = aba$ (PSD), $\det(s) = \prod s_i$ (nonneg).
Note: conex stores $W = w^{-1}$ internally; all formulas here use $w$ directly.}
\label{tab:dictionary}
\end{table}

\subsection*{Centering RHS verification}

For the symmetric cone with $F(s) = -\log\det(s)$:
\[
  \nabla^2 F(w) w - \nabla F(w) = P(w^{-1}) w + w^{-1} = w^{-1} + w^{-1} = 2w^{-1}.
\]
This gives $A^\top \cdot 2w^{-1}$ as the centering RHS.  In conex's convention
($W = w^{-1}$), this is $2A^\top W$. \checkmark

\section{Exponential cone: barrier derivatives}

The exponential cone is
$K_{\exp} = \operatorname{cl}\{(x,y,z) : y e^{x/y} \leq z,\, y > 0\}$,
with the 2-LHSCB ($\nu = 2$):
\[
  F(x,y,z) = -\log(z - y e^{x/y}) - \log y.
\]
Define the shorthand $e = e^{x/y}$, $s = x/y$, $u = z - ye > 0$ (the slack).

\subsection{Gradient}
\begin{align}
  F_x &= \frac{e}{u}, \qquad
  F_y = -\frac{e(s-1)}{u} - \frac{1}{y}, \qquad
  F_z = -\frac{1}{u}.
\end{align}

\subsection{Hessian}
Writing $H_{ij} = \partial^2 F / \partial x_i \partial x_j$ with
$(x_0, x_1, x_2) = (x, y, z)$:
\begin{align}
  H_{00} &= \frac{e(ey + u)}{u^2 y}, \\
  H_{01} &= -\frac{e\bigl(ey(s{-}1) + su\bigr)}{u^2 y}, \qquad
  H_{02} = -\frac{e}{u^2}, \\
  H_{11} &= \frac{e^2(s{-}1)^2}{u^2} + \frac{es^2}{uy} + \frac{1}{y^2}, \\
  H_{12} &= \frac{e(s{-}1)}{u^2}, \qquad
  H_{22} = \frac{1}{u^2}.
\end{align}

\subsection{Third derivatives (Christoffel symbols)}
\label{sec:christoffel}

The geodesic equation on the Hessian manifold is
\begin{equation}\label{eq:geodesic-ode}
  \ddot\gamma^\ell + \Gamma^\ell_{ij} \dot\gamma^i \dot\gamma^j = 0,
  \qquad
  \Gamma^\ell_{ij} = \frac{1}{2} H^{\ell m} F_{ijm},
\end{equation}
where $F_{ij\ell} = \partial^3 F / \partial x_i \partial x_j \partial x_\ell$
and $H^{\ell m}$ is the inverse Hessian.  The geodesic acceleration is
\[
  \ddot\gamma = -\tfrac{1}{2} H^{-1} \, T(v, v),
  \qquad
  T(v,v)_\ell = v^\top \frac{\partial H}{\partial x_\ell} v = \sum_{ij} F_{ij\ell} \, v_i v_j.
\]
The contraction $T(v,v)$ requires the 10 unique third derivatives of $F$
(symmetric in all three indices).  For the exponential cone:
\begin{align}
  F_{000} &= \frac{2e^3}{u^3} + \frac{3e^2}{u^2 y} + \frac{e}{u y^2}, \\
  F_{001} &= -\frac{e\bigl(2e^2 y^2(s{-}1) + euy(3s{-}1) + u^2(s{+}1)\bigr)}{u^3 y^2}, \\
  F_{002} &= -\frac{e(2ey + u)}{u^3 y}, \\
  F_{011} &= \frac{e\bigl(2esu y(s{-}1) + ey(2ey(s{-}1)^2 + s^2 u) + su^2(s{+}2)\bigr)}{u^3 y^2}, \\
  F_{012} &= \frac{e(2ey(s{-}1) + su)}{u^3 y}, \quad
  F_{022} = \frac{2e}{u^3}, \\
  F_{111} &= \frac{-2e^3 y^3(s{-}1)^3 + 3e^2 s^2 uy^2(1{-}s) - es^2 u^2 y(s{+}3) - 2u^3}{u^3 y^3}, \\
  F_{112} &= -\frac{e(2ey(s{-}1)^2 + s^2 u)}{u^3 y}, \quad
  F_{122} = \frac{2e(1{-}s)}{u^3}, \quad
  F_{222} = -\frac{2}{u^3}.
\end{align}
Evaluating $T(v,v)$ costs $3 \times 6 = 18$ multiply-adds (exploiting symmetry).
The full geodesic acceleration (including the $H^{-1}$ solve) is $O(1)$ per
integration step since $H$ is $3 \times 3$.

\subsection{St\"ormer--Verlet integrator}

Given position $\gamma(0) = s$ and velocity $\dot\gamma(0) = \alpha \cdot d$, the
geodesic is integrated by the symplectic St\"ormer--Verlet method with $N$ substeps
of size $\Delta t = 1/N$:
\begin{algorithmic}[1]
\For{$n = 1, \ldots, N$}
  \State $a \gets -\frac{1}{2} H(\gamma)^{-1} T(\dot\gamma, \dot\gamma)$
    \hfill (acceleration from Christoffel symbols)
  \State $\dot\gamma \gets \dot\gamma + \frac{\Delta t}{2} \, a$
    \hfill (half-step velocity)
  \State $\gamma \gets \gamma + \Delta t \, \dot\gamma$
    \hfill (full-step position)
  \State recompute $a$ at new $\gamma$
  \State $\dot\gamma \gets \dot\gamma + \frac{\Delta t}{2} \, a$
    \hfill (half-step velocity)
\EndFor
\end{algorithmic}
Each substep requires: one $3\times 3$ Hessian evaluation, one $T(v,v)$ evaluation
(10 third derivatives contracted with $v$), and one $3\times 3$ linear solve.
Total cost for $N = 8$ steps: $\sim$1000 flops.

\section{Bregman midpoint: geodesic approximation without third derivatives}
\label{sec:bregman}

The St\"ormer--Verlet integrator requires the Christoffel symbols $F_{ij\ell}$
(third derivatives of the barrier).  For low-dimensional cones (e.g., the
3D exponential cone) these are cheap.  For higher-dimensional or more complex
barriers, third derivatives may be expensive or unavailable.

The \emph{Bregman midpoint method} provides a second-order approximation to the
Levi-Civita geodesic using only the barrier gradient $\nabla F$ and Hessian $\nabla^2 F$.

\subsection{Dual flat structure of Hessian manifolds}

A Hessian manifold $(\operatorname{int}(K), \nabla^2 F)$ carries two flat
connections:
\begin{itemize}
  \item \emph{Primal flat ($\nabla$-geodesics):} straight lines $\gamma(t) = s + tv$
    in the original coordinates $s$.
  \item \emph{Dual flat ($\nabla^*$-geodesics):} straight lines
    $\lambda(t) = \lambda_0 + tw$ in the dual coordinates $\lambda = -\nabla F(s)$.
\end{itemize}
The Levi-Civita geodesic (which preserves the Riemannian distance and stays in
$\operatorname{int}(K)$) lies ``between'' these two flat geodesics, agreeing with
both to first order.

\subsection{The midpoint construction}

Given starting point $s$ and direction $d$, the Bregman midpoint step computes
$s_{\text{new}} \approx \Exp_s(\alpha \cdot d)$ as follows:
\begin{enumerate}
  \item \textbf{Primal half-step:}
    $s_{1/2} = s + \frac{\alpha}{2} d$ \quad (straight line in primal coordinates).

  \item \textbf{Map to dual at midpoint:}
    $\lambda_{1/2} = -\nabla F(s_{1/2})$.

  \item \textbf{Map to dual at start:}
    $\lambda_0 = -\nabla F(s)$.

  \item \textbf{Dual extrapolation} (straight line in dual coordinates):
    $\lambda_1 = 2\lambda_{1/2} - \lambda_0$.

    This is the midpoint rule: the dual velocity
    $\dot\lambda \approx (\lambda_{1/2} - \lambda_0) / (\alpha/2)$, so
    $\lambda_1 = \lambda_0 + \alpha \dot\lambda = 2\lambda_{1/2} - \lambda_0$.

  \item \textbf{Invert the gradient map:} find $s_{\text{new}}$ such that
    $-\nabla F(s_{\text{new}}) = \lambda_1$.  This requires solving a
    nonlinear system (Newton's method on the $m$-dimensional equation
    $\nabla F(s) + \lambda_1 = 0$).
\end{enumerate}

\subsection{Computational cost}

\begin{center}
\begin{tabular}{@{}lcc@{}}
\toprule
\textbf{Operation} & \textbf{Evaluations needed} & \textbf{Cost (3D cone)} \\
\midrule
Step 1 (primal half-step) & --- & $O(1)$ \\
Steps 2--3 (gradient at two points) & $2 \times \nabla F$ & $O(1)$ \\
Step 4 (dual extrapolation) & --- & $O(1)$ \\
Step 5 (gradient inversion) & $3$--$5 \times (\nabla F + \nabla^2 F)$ & $O(1)$ \\
\midrule
\textbf{Total} & $\sim$5 gradient + 5 Hessian & $\sim$200 flops \\
\bottomrule
\end{tabular}
\end{center}

No third derivatives are required.  The dominant cost is the Newton solve
in step~5 (inverting the gradient map), which for the 3D exponential cone
converges in 3--5 iterations of a $3 \times 3$ system.

\subsection{Accuracy}

The Bregman midpoint is \emph{second-order accurate}: the error is $O(\alpha^3)$
compared to the true Levi-Civita geodesic.  At $\alpha = 0.1$, the Bregman
midpoint is $\sim$200$\times$ more accurate than the Euler step (which is
first-order, $O(\alpha^2)$ error).

For comparison, the St\"ormer--Verlet integrator with $N = 8$ substeps achieves
higher accuracy ($\sim$fourth order) but requires third derivatives.

\subsection{When to use each method}

\begin{center}
\begin{tabular}{@{}lll@{}}
\toprule
& \textbf{St\"ormer--Verlet} & \textbf{Bregman midpoint} \\
\midrule
Derivatives needed & $\nabla F$, $\nabla^2 F$, $\partial^3 F$ & $\nabla F$, $\nabla^2 F$ only \\
Accuracy & High (symplectic, $\sim$4th order) & Second order \\
Cost per step (3D) & $\sim$1000 flops & $\sim$200 flops \\
Aggressive steps & Yes ($\alpha$ up to 1) & Conservative ($\alpha \ll 1$) \\
Best for & Low-dim cones with known $\partial^3 F$ & General barriers \\
\bottomrule
\end{tabular}
\end{center}

For the exponential cone ($m = 3$), the third derivatives are closed-form
(Section~\ref{sec:christoffel}) and the St\"ormer--Verlet integrator is preferred.
For general barriers where $\partial^3 F$ is unavailable or expensive, the Bregman
midpoint provides a practical alternative that still preserves the cone interior
and achieves second-order accuracy.

\subsection{Dual geodesic property}

The gradient map $s \mapsto -\nabla F(s)$ is an isometry between
$(\operatorname{int}(K), \nabla^2 F)$ and $(\operatorname{int}(K^*), \nabla^2 F^*)$.
Therefore, if $\gamma(t)$ is a geodesic in $K$, then $\lambda(t) = -\nabla F(\gamma(t))$
is a geodesic in $K^*$.  The Bregman midpoint exploits this by stepping in dual
coordinates (where the geodesic is approximately a straight line) and mapping back.

This property was verified numerically for the exponential cone: the dual
energy $|\dot\lambda|^2_{H^{-1}}$ is constant to within 0.3\% along the
integrated geodesic.

\section{Exponential cone: numerical results}

For a product of $m$ exponential cones with $p$ variables, the Gram is
\[
  G = \sum_{i=1}^m A_i^\top \nabla^2 F_i(w_i) A_i,
\]
and the geodesic step is taken component-wise on each 3D cone.

\paragraph{Centered initial point.}
Choose $b_i \in \operatorname{int}(K_{\exp})$ and set
$c = -\sum_i A_i^\top \nabla F_i(b_i)$, ensuring $x = 0$ is on the central path
at $\mu = 1$.  This is the analogue of $c = A^\top e$ for symmetric cones.

\paragraph{Results.}
On a test problem with $m = 3$ cones and $p = 4$ variables,
the algorithm converges in 4 factorizations with $k$ growing $\sim$10$\times$
per step and duality gap decreasing $\sim$100$\times$ per step.  The dual
residual $\|A^\top \lambda - c\|$ is $\sim 10^{-16}$ (machine zero) at each
iteration.

\end{document}
